\documentclass[10pt,letterpaper]{article}
\usepackage[T1]{fontenc}
\usepackage{amssymb}
\usepackage{amsmath}
\usepackage{enumerate}
\usepackage[margin=0.5in]{geometry}
\usepackage{xcolor}

\definecolor{TangGreen}{HTML}{377459}

\title{Design and Analysis of Algorithms Homework 2}
\author{Jeremy Chen}
\begin{document}
	\maketitle

\noindent\textbf{Problem 1. \hfill (20 points)} 

\noindent\textbf{The greedy set algorithm is not better than} $\log n$. Recall the set cover problem. We are given a universe $U$ of $n$ elements and a collection $\mathcal{S} = \{ S_{1}, \dots, S_{m} \}$ of subsets of $U$ whose union is $U$. We must choose as few of these subsets as possible so that their union is all of $U$. The greedy algorithm repeatedly selects a set covering the largest number of still-uncovered elements, until nothing remains uncovered.

In lecture we proved that greedy uses at most $k\ln n$ sets, where $k$ is the number of sets in an optimal cover. This problem shows that the analysis cannot be improved beyond a constant factor. 

\begin{enumerate}[(a)]

\item \textbf{(5 points)} Give a mall explicit instance on which greedy does not return an optimal cover. State the sets greedy selects, the sets an optimal cover selects, and both sizes.

\color{TangGreen}

\color{black}

\item \textbf{(11 points)} For each integer $k \geq 1$, construct an instance with $n = 2^{k + 1} - 2$ elements in which 

\begin{itemize}

\item some cover uses only 2 sets, and

\item the greedy algorithm uses $k$-sets.

\end{itemize}

Describe the universe and collection precisely, then prove both claims. You may assume greedy breaks ties arbitrarily; your construction should not depend on how. 

\color{TangGreen}
Start by resolving the the base case, where $k = 1$. For this case, there will be $2^{1 + 1} - 2 = 2$ elements. Therefore, the universe is
$$U = \{ \ell_{1}, \ell_{2} \}$$
and the collections are 
$$C_{1} = \{ \ell_{1} \}, C_{2} = \{ \ell_{2} \}, C_{3} = \{ \ell_{1}, \ell_{2} \}$$
Greedy will select the set with the largest number of uncovered elements, $C_{3}$, and use $k = 1$ sets. Some cover could use two sets $C_{1}$ and $C_{2}$. For another case, where $k = 2$, there are $2^{2 + 1} - 1 = 6$ elements. The universe is 
$$U = \{ \ell_{1}, \ell_{2}, \dots, \ell_{6} \}$$
and the collection is then defined as 
$$C_{1} = \{ \ell{1}, \ell_{3}, \ell_{5} \}, C_{2} = \{ \ell_{2}, \ell_{4}, \ell_{6} \}, C_{3} = \{ \ell_{1}, \ell_{2}, \ell_{3}, \ell_{4} \}, C_{4} = \{ \ell_{5}, \ell_{6} \}$$
At this point, a pattern emerges. A universe would contain $2^{k + 1} - 2$ unique elements, and the collection consists of two sets with $2^{k} - 1$ elements each, and another $k$ disjunct sets, with $2^{k}$, $2^{k - 1}$, and so on elements per.

At the start, the greedy algorithm would select the set with $2^{k}$, since it has more uncovered elements than the sets with $2^{k} - 1$ elements. 
\color{black}

\end{enumerate}

\end{document}